\uuid{V60M}
\titre{Matrice d'une application linéaire/Changement de bases/Puissance de matrices}
\auteur{hanani ; mzari}
\organisation{exo7}
\licence{CC BY-NC-SA 4.0}
\datecreate{2020-01-20}
\niveau{L1}
\module{Algèbre}
\chapitre{Matrice}
\sousChapitre{Matrice et application linéaire}
\difficulte{2}
\theme{Matrice d'une application linéaire, Changement de bases, Puissance de matrices}
\exo7qcmid{lille-2.matrice-application.014}

\contenu{
\begin{question}
\qclassification{108.03}
\qtags{motcle=Matrice d'une application linéaire/Changement de bases/Puissance de matrices}
Soit $f$ l'endomorphisme de $\Rr^2$ dont la matrice dans la base canonique ${\cal B}=\{e_1,e_2\}$ est :
$$A=\left(\begin{array}{rc}1&3\\
1&-1\\ \end{array}\right).$$
Soit ${\cal B}' = \{ u_1, u_2\}$, où  $u_1=(3,1), u_2=(1,-1)$, une base de $\Rr^2$. On note $P$ la matrice de passage de la base ${\cal B}$ à la base ${\cal B}'$ et $B$ La matrice de $f$ dans la base ${\cal B}'$.
\vskip1mm
Définition : Soit $E$ un espace vectoriel muni de deux bases ${\cal {B}}$ et ${\cal {B}}'$. La matrice de passage de la base ${\cal {B}}$ à la base  ${\cal {B}}'$ est la matrice de l'identité de $E$ de la base ${\cal {B}}'$ à la base  ${\cal {B}}$. Autrement dit, c'est la matrice dont la jième colonne est constituée des coordonnées du jième vecteur de la base ${\cal {B}'}$ dans la base  ${\cal {B}}$.
\vskip1mm
Quelles sont les assertions vraies ?
\begin{answers}
\good{$P=\left(\begin{array}{rc}3&1\\1&-1\\
\end{array}\right)  $}.
\bad{$P^{-1}= \left(\begin{array}{rc}3&1\\1&-1\\
\end{array}\right).$}
\bad{$B= \left(\begin{array}{rc}-2&0\\0&2\\
\end{array}\right)  $}.
\good{$A^n= 2^{n-2}\left(\begin{array}{rc}
3+(-1)^n&3-3(-1)^n\\
 1-(-1)^n &1+3(-1)^n\\
\end{array}\right) $, pour tout entier $n\ge1$}.
\end{answers}
\begin{explanations} Soit $E$ un espace vectoriel de dimension finie, muni de deux bases ${\cal {B}}$ et ${\cal {B}}'$
et  $P$ la matrice de passage de la base ${\cal {B}}$ à la base ${\cal {B}}'$. Soit $f$ un endomorphisme de $E$ de matrice
$A$ (resp. $B$) dans la base ${\cal {B}}$ (resp. ${\cal {B}}'$).
Alors, on a la relation : $AP=PB$. De cette relation, on déduit que $A^n=PB^nP^{-1}$. $P^{-1}= \frac{1}{4}\left(\begin{array}{rc}
1&1\\ 1&-3\\ \end{array}\right)$,  $B= \left(\begin{array}{rc}
2&0\\0&-2\\ \end{array}\right)$ et que
$$A^n= 2^{n-2}\left(\begin{array}{rc}3+(-1)^n&3-3(-1)^n\\
1-(-1)^n &1+3(-1)^n\end{array}\right),\mbox{  pour tout entier }n\ge1.$$
\end{explanations}
\end{question}
}
